%%%%%%%%%%%%%%%%%%%%%%%%%%%%%%%%%%%%%%%%
%
% o CRONOGRAMA é obrigatório 
% APENAS PARA PRÉ-PROJETO de PFC
%
%%%%%%%%%%%%%%%%%%%%%%%%%%%%%%%%%%%%%%%%

\chapter{CRONOGRAMA}

O cronograma deste trabalho está organizado em fases que abrangem desde a pesquisa teórica até a implementação prática do projeto. A primeira etapa consiste no levantamento bibliográfico e estudo dos fundamentos teóricos necessários para o desenvolvimento do sistema. Em seguida, serão realizados o projeto e a simulação do circuito, seguidos pela montagem e testes do protótipo. A fase final inclui a análise dos resultados, a redação do trabalho e a preparação para a apresentação. Cada etapa foi planejada de forma a garantir o cumprimento dos prazos e a qualidade do trabalho entregue, conforme tabela~\ref{tab:Cronograma}.

\begin{table}[htb]
  \centering
  \caption{Cronograma de Execução do Projeto Final de Curso}
  \begin{tabular}{m{6cm}cccccc} % l=Left(textos); r=Right(números); o comando m{6cm} fixa a primeira coluna em 6cm, faz quebra de linha e centraliza as demais colunas na vertical
    \Linha % linha grossa
    {\bf CRONOGRAMA PFC-2025} & {\bf Jul} & {\bf Ago}  & {\bf Set} & {\bf Out} & {\bf Nov} & {\bf Dez}  \\
    \Linha % linha grossa      
    Levantamento bibliográfico & X &   &   &   &   &     \\
    \hline % linha fina
    Alinhamento da proposta    & X &   &   &   &   &     \\
    \hline % linha fina
    Projeto                    &   & X &   &   &   &     \\
    \hline % linha fina
    Montagem e experimentos    &   & X & X &   &   &     \\
    \hline % linha fina
    Levantamento de dados,   
    criação de tabelas e 
    gráficos                   &   &   &  X &   &   &    \\
    \hline
    Escrita da monografia      &   &   &  X &  X  &   &  \\
    \hline % linha fina
    Revisão da monografia      &   &   &  X &  X  &   &  \\
    \hline  
    Defesa perante Banca       &   &   &    &     &  X &    \\
    \Linha % linha grossa
  \end{tabular}    
  \FonteTab{Próprio Autor}
  \label{tab:Cronograma}
\end{table}
